\begin{nomenclature}
\zihao{5}
\renewcommand\arraystretch{1.15}

\begin{longtable}{p{3.2cm}p{10.8cm}}
  \toprule
  符号 & 代表意义 \\
  \midrule
  \endfirsthead

  \toprule
  符号 & 代表意义 \\
  \midrule
  \endhead

  \bottomrule
  \endfoot

  \bottomrule
  \endlastfoot

  \(s_u(t)\) & 目标发射的等效基带信号 \\
  \(\nu_l(t)\) & 第 \(l\) 条传播路径对应的多普勒频移 \\
  \(\sigma_{\tau}\) & 均方根时延扩展，用于表征多径传播复杂度 \\
  \(\mathbf{p}_{k}^{\mathrm{RSSI}}\) & 第 \(k\) 个样本由 RSSI 分支得到的二维位置估计 \\
  \(\mathbf{p}_{k}^{\mathrm{AoA}}\) & 第 \(k\) 个样本由 AoA 分支得到的二维位置估计 \\
  \(\mathbf{p}_{k}^{\mathrm{PDR}}\) & 第 \(k\) 个样本由 PDR 递推分支得到的二维位置估计 \\
  \(\hat{\mathbf{p}}_{k}^{\mathrm{Fused}}\) & 第 \(k\) 个样本的融合定位结果 \\
  \(\mathbf{p}_{k}^{\mathrm{th}}\) & 第 \(k\) 个样本的真实位置 \\
  \(\mathbf{P}_k\) & 第 \(k\) 个样本的三路定位结果矩阵，\(\mathbf{P}_k\in\mathbb{R}^{2\times3}\) \\
  \(\mathbf{t}_k\) & 第 \(k\) 个样本的输入特征向量 \\
  \(\mathbf{q}_k^\ast\) & 离线反解得到的最优融合权重标签 \\
  \(\hat{\mathbf{q}}_k\) & XGBoost 预测得到的融合权重向量 \\
  \(\mathbf{T},\mathbf{Q}^{\ast}\) & 样本特征矩阵与离线反解权重标签矩阵 \\
  \(\sigma_{\mathrm{RSSI}}^2\) & RSSI 在滑动窗口内的方差，用于表征观测波动性 \\
  \(R_q\) & 权重变化率，用于衡量相邻样本融合权重波动程度 \\
  \(X_{i,j}\) & 第 \(i\) 架无人机对第 \(j\) 个目标样本的单节点观测特征向量 \\
  \(N\) & 协同定位网络中的无人机节点数量 \\
  \(M\) & 目标区域离散后的栅格类别数 \\
  \(\rho_{\mathrm{cov}}\) & 多无人机协同观测中的有效节点覆盖率 \\
  \(\mathbf{h}_i\) & 第 \(i\) 个节点经过图注意力网络提取后的高层特征表示 \\
  \(\alpha_{ij}\) & 图注意力机制中节点 \(j\) 对节点 \(i\) 的注意力权重 \\
  \(\mathbf{z}^{T},\mathbf{z}^{S}\) & 教师网络与学生网络输出的分类得分向量 \\
  \(T_d\) & 知识蒸馏中的温度系数 \\
  \(\alpha,\beta\) & 交叉熵损失与蒸馏损失的权重系数 \\
  \(\mathbf{p}^{\mathrm{gt}}\) & 目标真实位置向量 \\
  \(\hat{\mathbf{p}}\) & 模型预测得到的目标位置向量 \\
  \(e\) & 单样本定位误差 \\
\end{longtable}

\end{nomenclature}
